\section{刑罚的执行和消灭}
\begin{enumerate}
  \item 可参考刑诉。
  \item 减刑、\kd{假释}。考的多。
  \item 漏罪、新罪问题。\zd{是否可以撤销、是否构成累犯、是否可以再次适用}。
        \begin{enumerate}
          \item \zd{期后}发现\zd{漏罪}都只能另行\zd{起诉}，其他情形都应该撤销。
          \item 除了假释期后发现\zd{新罪}有可能构成\zd{累犯}，其他都不构成。
          \item 除了\zd{缓刑}发现\zd{新罪}可以再次适用\zd{假释}，其他都不得再次适用缓刑/假释。
          \item 刑期计算规则不变。
        \end{enumerate}
  \item \kd{追诉时效}。
        \begin{enumerate}
          \item 最高刑不满。
                \begin{enumerate}
                  \item 5 年 \yjt\, 5 年。
                  \item 5 - 10 \yjt\, 10。
                  \item 10+ \yjt\, 15。
                  \item 无期、死刑 \yjt\, 20。仅这一档 > 20 特别追诉可以找最高检批。
                  \item \al{最高刑 10+，过了 15 年，特别追诉可以找最高检 \cuo}。
                \end{enumerate}
          \item 计算。其实是起算日和截止日的问题。
                \begin{enumerate}
                  \item 起算日。
                        \begin{enumerate}
                          \item 故意犯罪 \yjt 产生危险就是\zd{成立}，实害结果是既遂。\al{1/10 挪用公款，7/10 归还，4/11 起算 \dui}。
                          \item 过失犯罪 \yjt 实害结果发生之日。
                          \item 连续犯、继续犯 \yjt \zd{犯罪行为终了之日}。
                        \end{enumerate}
                \end{enumerate}
          \item 延长。\zd{前提是在时效内}。
                \begin{enumerate}
                  \item 公检法已立案侦查 + 犯罪嫌疑人逃避侦查 \yjt 无限延长。
                  \item 被害人提出控告 + 不立案 \yjt 无限延长。\al{孙小果}。
                \end{enumerate}
          \item 中断（清零重计算）。追诉期限内犯了新罪，原来的罪的追诉时效统统作废清零。
          \item \zd{共犯。分别计算，互不影响}。
        \end{enumerate}
\end{enumerate}
